\documentclass[
    language=english,
    doctype=risk-assessment,
    institution=none,
]{../../omnilatex}

\RequirePackage{config/document-settings}

\begin{document}

\maketitle

\section*{Executive Summary}

This risk assessment evaluates threats to NovaPulse Technologies' platform
operations, data security, regulatory compliance, and strategic objectives
as of Q2 2026. The assessment covers 18 identified risks across five
categories, evaluated using a 5$\times$5 likelihood-impact matrix aligned
with ISO 31000 and NIST SP 800-30 frameworks.

Key findings:

\begin{itemize}
    \item 3 risks rated \textbf{Critical} (score 15--25): cybersecurity
          breach, key talent attrition in ML engineering, and data residency
          compliance failure.
    \item 5 risks rated \textbf{High} (score 10--14): cloud infrastructure
          outage, competitive pricing pressure, regulatory change (EU AI Act),
          third-party vendor failure, and intellectual property dispute.
    \item 6 risks rated \textbf{Medium} (score 5--9) and 4 risks rated
          \textbf{Low} (score 1--4).
    \item Aggregate risk exposure decreased 12\% from Q1 2026 due to
          completion of SOC 2 Type II audit and launch of employee retention
          program.
\end{itemize}

\begin{table}[h]
    \centering
    \begin{tabular}{lrr}
        \toprule
        \textbf{Summary Metric} & \textbf{Value} \\
        \midrule
        Total Risks Identified       & 18 \\
        Critical Risks               & 3  \\
        High Risks                   & 5  \\
        Medium Risks                 & 6  \\
        Low Risks                    & 4  \\
        Risks with Mitigation Plans  & 18 (100\%) \\
        Risks Requiring Escalation   & 4  \\
        Assessment Date              & June 15, 2026 \\
        Next Assessment Due          & September 15, 2026 \\
        \bottomrule
    \end{tabular}
    \caption{Risk assessment summary.}
\end{table}

\section{Methodology}

\subsection{Risk Rating Framework}

Each risk is evaluated on two dimensions:

\begin{description}
    \item[Likelihood.] The probability of the risk event occurring within
        the next 12 months, rated 1 (Very Low, $<$5\%) through 5 (Almost
        Certain, $>$80\%).
    \item[Impact.] The severity of consequences if the risk materializes,
        rated 1 (Negligible) through 5 (Catastrophic), considering
        financial, operational, reputational, and regulatory dimensions.
\end{description}

Risk score = Likelihood $\times$ Impact. Scores are mapped to risk levels:

\begin{table}[h]
    \centering
    \begin{tabular}{llc}
        \toprule
        \textbf{Risk Level} & \textbf{Score Range} & \textbf{Response} \\
        \midrule
        Critical  & 15--25 & Immediate action required; escalate to C-suite \\
        High      & 10--14 & Active mitigation plan; monthly review \\
        Medium    & 5--9   & Monitor with defined triggers; quarterly review \\
        Low       & 1--4   & Accept and monitor; annual review \\
        \bottomrule
    \end{tabular}
    \caption{Risk level classification.}
\end{table}

\subsection{Assessment Scope}

This assessment covers:

\begin{enumerate}
    \item Technology and cybersecurity risks.
    \item Operational and business continuity risks.
    \item People and organizational risks.
    \item Regulatory and compliance risks.
    \item Market and competitive risks.
\end{enumerate}

\section{Risk Matrix}

\begin{figure}[h]
    \centering
    \begin{tikzpicture}[
        cell/.style={minimum width=2.4cm,minimum height=1.2cm,font=\small},
    ]
        \matrix (risk) [matrix of nodes,
            nodes={cell,anchor=center},
            column sep=0pt, row sep=0pt,
            column 1/.style={nodes={draw=none,minimum width=2.2cm,fill=none}},
            row 1/.style={nodes={draw=none,fill=none}},
            column 2/.style={nodes={fill=green!15,draw=gray!50}},
            column 3/.style={nodes={fill=green!30,draw=gray!50}},
            column 4/.style={nodes={fill=yellow!35,draw=gray!50}},
            column 5/.style={nodes={fill=orange!35,draw=gray!50}},
            column 6/.style={nodes={fill=red!30,draw=gray!50}},
            row 2/.style={nodes={fill=red!15,draw=gray!50}},
            row 3/.style={nodes={fill=orange!15,draw=gray!50}},
            row 4/.style={nodes={fill=yellow!20,draw=gray!50}},
            row 5/.style={nodes={fill=green!20,draw=gray!50}},
            row 6/.style={nodes={fill=green!10,draw=gray!50}},
        ] {
            & \textbf{Negligible} & \textbf{Minor} & \textbf{Moderate} & \textbf{Major} & \textbf{Catastrophic} \\
            \textbf{Almost Certain} & & & R5 & & \\
            \textbf{Likely}   & & & R11 & R1,R2 & R3 \\
            \textbf{Possible} & R17 & R14 & R8,R12 & R6,R7 & \\
            \textbf{Unlikely} & R18 & R13,R16 & R9,R10 & R4 & \\
            \textbf{Rare}     & & R15 & & & \\
        };
        \node[above=0.3cm of risk-1-3,font=\small\bfseries] {Impact $\rightarrow$};
        \node[left=0.1cm of risk-3-1,rotate=90,font=\small\bfseries,anchor=south] {Likelihood $\rightarrow$};
    \end{tikzpicture}
    \caption{Risk matrix: likelihood vs.\ impact with identified risks plotted.}
\end{figure}

\section{Detailed Risk Register}

\subsection{Technology \& Cybersecurity Risks}

\begin{table}[h]
    \centering
    \small
    \begin{tabular}{llcccl}
        \toprule
        \textbf{ID} & \textbf{Risk} & \textbf{L} & \textbf{I} & \textbf{Score} & \textbf{Mitigation Strategy} \\
        \midrule
        R1 & Cybersecurity breach / data exfiltration & 4 & 4 & 16 &
            SOC 2 Type II certified; quarterly pen testing; 24/7 SOC
            monitoring; zero-trust architecture; \$5M cyber insurance \\
        R2 & Ransomware attack on production systems & 4 & 4 & 16 &
            Immutable backups with 30-day retention; network segmentation;
            endpoint detection (CrowdStrike); incident response plan tested
            quarterly \\
        R3 & ML model poisoning / adversarial attack & 4 & 5 & 20 &
            Input validation; model integrity monitoring; adversarial
            robustness testing; human-in-the-loop validation for critical
            predictions \\
        \bottomrule
    \end{tabular}
    \caption{Technology and cybersecurity risk details.}
\end{table}

\subsection{Operational Risks}

\begin{table}[h]
    \centering
    \small
    \begin{tabular}{llcccl}
        \toprule
        \textbf{ID} & \textbf{Risk} & \textbf{L} & \textbf{I} & \textbf{Score} & \textbf{Mitigation Strategy} \\
        \midrule
        R4 & Cloud provider outage (AWS regional) & 2 & 4 & 8 &
            Multi-AZ deployment; active-active failover to EU region
            (Project Atlas); DR testing quarterly; 99.99\% SLA target \\
        R5 & Third-party vendor failure (key supplier) & 3 & 3 & 9 &
            Vendor risk assessments; contractual SLAs with penalties;
            identified backup vendors for critical services; 90-day buffer
            stock for hardware \\
        R6 & Insider threat / unauthorized access & 3 & 4 & 12 &
            Role-based access control; least-privilege principles; audit
            logging; background checks; mandatory security training;
            anonymous reporting channel \\
        R7 & Data loss / corruption & 3 & 4 & 12 &
            Automated backups every 15 minutes; cross-region replication;
            checksums and integrity validation; tested restore procedures
            monthly \\
        \bottomrule
    \end{tabular}
    \caption{Operational risk details.}
\end{table}

\subsection{People Risks}

\begin{table}[h]
    \centering
    \small
    \begin{tabular}{llcccl}
        \toprule
        \textbf{ID} & \textbf{Risk} & \textbf{L} & \textbf{I} & \textbf{Score} & \textbf{Mitigation Strategy} \\
        \midrule
        R8 & Key talent attrition (ML engineers) & 4 & 4 & 16 &
            Competitive compensation (80th percentile); retention bonuses
            for critical roles; expanded equity grants; EAP enhancements;
            career development programs \\
        R9 & Knowledge concentration risk & 3 & 3 & 9 &
            Cross-training programs; documented architecture decisions;
            pair programming rotation; minimum 2 engineers per critical
            system \\
        R10 & Onboarding bottleneck (new hires) & 3 & 3 & 9 &
            Structured 30-60-90 onboarding program; buddy system;
            documented runbooks; automated environment setup \\
        \bottomrule
    \end{tabular}
    \caption{People risk details.}
\end{table}

\subsection{Regulatory \& Compliance Risks}

\begin{table}[h]
    \centering
    \small
    \begin{tabular}{llcccl}
        \toprule
        \textbf{ID} & \textbf{Risk} & \textbf{L} & \textbf{I} & \textbf{Score} & \textbf{Mitigation Strategy} \\
        \midrule
        R11 & EU AI Act non-compliance & 3 & 3 & 9 &
            Dedicated compliance team; privacy-by-design architecture;
            automated model documentation; third-party audit scheduled
            Q4 2026 \\
        R12 & Data residency violation & 3 & 4 & 12 &
            Project Atlas delivers configurable data residency; automated
            compliance monitoring; legal review of all data flows; customer
            data classification system \\
        R13 & GDPR enforcement action & 2 & 3 & 6 &
            Appointed DPO; DPIAs conducted for all new features; consent
            management platform; data subject request automation \\
        \bottomrule
    \end{tabular}
    \caption{Regulatory and compliance risk details.}
\end{table}

\subsection{Market \& Strategic Risks}

\begin{table}[h]
    \centering
    \small
    \begin{tabular}{llcccl}
        \toprule
        \textbf{ID} & \textbf{Risk} & \textbf{L} & \textbf{I} & \textbf{Score} & \textbf{Mitigation Strategy} \\
        \midrule
        R14 & Competitive pricing pressure & 3 & 2 & 6 &
            Differentiate on SLA, compliance, and edge capabilities;
            value-based selling; flexible pricing tiers \\
        R15 & Macroeconomic downturn reducing AI spending & 2 & 2 & 4 &
            Position as cost-reduction tool; multi-year contracts;
            diversified customer base across 8 industries \\
        R16 & Open-source alternative commoditizes core features & 2 & 3 & 6 &
            Invest in proprietary ML models; domain-specific solutions;
            managed service layer; customer lock-in through integrations \\
        R17 & Intellectual property dispute & 3 & 1 & 3 &
            IP insurance; regular freedom-to-operate analyses; patent
            portfolio review \\
        R18 & Reputational damage from customer incident & 2 & 1 & 2 &
            Incident communication plan; customer notification procedures;
            transparent post-mortem process; cyber insurance \\
        \bottomrule
    \end{tabular}
    \caption{Market and strategic risk details.}
\end{table}

\section{Mitigation Effectiveness}

\begin{table}[h]
    \centering
    \small
    \begin{tabular}{lccccc}
        \toprule
        \textbf{Risk} & \textbf{Inherent Score} & \textbf{Res.\ L} & \textbf{Res.\ I} & \textbf{Residual Score} & \textbf{Reduction} \\
        \midrule
        R3: ML model poisoning  & 20 & 3 & 4 & 12 & $-$40\% \\
        R1: Cybersecurity breach & 16 & 3 & 3 & 9  & $-$44\% \\
        R2: Ransomware attack    & 16 & 3 & 3 & 9  & $-$44\% \\
        R8: Key talent attrition & 16 & 3 & 3 & 9  & $-$44\% \\
        R6: Insider threat       & 12 & 2 & 3 & 6  & $-$50\% \\
        R7: Data loss            & 12 & 2 & 3 & 6  & $-$50\% \\
        R12: Data residency      & 12 & 2 & 3 & 6  & $-$50\% \\
        R4: Cloud outage         &  8 & 1 & 3 & 3  & $-$63\% \\
        R11: EU AI Act           &  9 & 2 & 2 & 4  & $-$56\% \\
        \bottomrule
    \end{tabular}
    \caption{Residual risk scores after mitigation (top 9 risks by score).}
\end{table}

\section{Risk Response Actions}

\begin{table}[h]
    \centering
    \small
    \begin{tabular}{llcl}
        \toprule
        \textbf{Risk} & \textbf{Action} & \textbf{Owner} & \textbf{Due Date} \\
        \midrule
        R3  & Complete adversarial robustness testing framework &
            Marcus Chen & Aug 2026 \\
        R8  & Finalize retention bonus program for 12 critical ML engineers &
            Priya Sharma & Jul 2026 \\
        R12 & Launch data residency beta in Project Atlas &
            Taylor Brooks & Oct 2026 \\
        R6  & Complete zero-trust migration for production environments &
            Casey Diaz & Sep 2026 \\
        R1  & Obtain SOC 2 Type II certification (already complete; annual
            renewal scheduled) &
            Riley Chen & Jun 2027 \\
        R2  & Conduct tabletop ransomware response exercise with full team &
            Jordan Mitchell & Aug 2026 \\
        \bottomrule
    \end{tabular}
    \caption{Key risk response actions with owners and deadlines.}
\end{table}

\section{Governance \& Review}

\begin{itemize}
    \item \textbf{Monthly risk review.} The Risk Committee (CTO, CISO,
        VP Engineering, Legal) meets monthly to review the top 10 risks,
        track mitigation progress, and escalate new threats.
    \item \textbf{Quarterly board reporting.} A risk dashboard is presented
        to the Board of Directors at each quarterly meeting, including
        changes in risk scores and status of critical mitigation actions.
    \item \textbf{Incident-triggered reassessment.} Any P0 or P1 incident
        triggers an immediate reassessment of the affected risk category.
    \item \textbf{Annual comprehensive review.} A full risk assessment is
        conducted annually with input from all department heads and external
        advisors.
\end{itemize}

\section{Signatures}

\begin{table}[h]
    \centering
    \begin{tabular}{lll}
        \toprule
        \textbf{Role} & \textbf{Name} & \textbf{Date} \\
        \midrule
        CTO / Risk Committee Chair & Marcus Chen & \underline{\hspace{3cm}} \\
        CISO                        & (TBD) & \underline{\hspace{3cm}} \\
        VP Engineering              & Priya Sharma & \underline{\hspace{3cm}} \\
        General Counsel             & (TBD) & \underline{\hspace{3cm}} \\
        \bottomrule
    \end{tabular}
    \caption{Assessment approval signatures.}
\end{table}

\printbibliography[heading=none]

\end{document}
